% Generated draft: requires mathematical and source review.
\documentclass[11pt]{article}
\usepackage[margin=25mm]{geometry}
\usepackage{fontspec}
\setmainfont{Latin Modern Roman}
\usepackage{amsmath,amssymb,mathtools,unicode-math}
\setmathfont{Latin Modern Math}
\usepackage{enumitem,needspace,xurl,hyperref}
\setlength{\parindent}{0pt}
\setlength{\parskip}{5pt}
\setlength{\emergencystretch}{3em}
\newcommand{\N}{\mathbb{N}}
\newcommand{\Z}{\mathbb{Z}}
\newcommand{\Q}{\mathbb{Q}}
\newcommand{\R}{\mathbb{R}}
\newcommand{\C}{\mathbb{C}}
\newcommand{\F}{\mathbb{F}}
\newcommand{\A}{\mathbb{A}}
\newcommand{\PP}{\mathbb{P}}
\newcommand{\E}{\mathbb{E}}
\newcommand{\eps}{\varepsilon}
\newcommand{\dd}{\,\mathrm{d}}
\begin{document}
\section*{UnsolvedMath ID: 30001639}
\textbf{Source problem code:} catalogue entry
\subsection{Definitions and mathematical statement}
The supplied record poses four related questions about compactifications of moduli spaces and thick parts of Teichmüller space. Several objects and comparison maps are unspecified, so the questions below retain those qualifications rather than assert a fully determined conjecture.

For the usual closed-surface interpretation, let \(S_g\) be a closed connected oriented surface of genus \(g\geq 2\), and let \(T(S_g)\) denote its Teichmüller space of marked conformal structures. Using the hyperbolic metric of curvature \(-1\) associated with such a structure \(X\), define
\[
\operatorname{sys}(X)=\inf_{\gamma}\ell_X(\gamma),
\]
where \(\gamma\) ranges over essential simple closed curves and \(\ell_X(\gamma)\) is the length of the geodesic representative. For \(\varepsilon>0\), distinguish
\[
T^{\geq\varepsilon}(S_g)=\{X\in T(S_g):\operatorname{sys}(X)\geq\varepsilon\},\qquad
T^{>\varepsilon}(S_g)=\{X\in T(S_g):\operatorname{sys}(X)>\varepsilon\}.
\]
These definitions explain conventional possible meanings of the source notation \(T^\varepsilon(S_g)\); the record does not specify the surface hypotheses, genus range, permitted values of \(\varepsilon\), metric normalization, or boundary convention. In particular, the closed-surface interpretation above is conditional, not an additional supplied hypothesis.

\begin{enumerate}
\item Which classes of the Deligne--Mumford compactification come from the Penner compactification? A precise version requires the moduli problem, the particular Penner compactification, the meaning of ``classes,'' their coefficient system and degrees, and the comparison map or correspondence through which they come from that compactification. None of these choices is supplied sufficiently to define an image of a specified map.
\item Is the Sullivan space homotopic to a suitably compactified moduli space? If ``homotopic'' means ``homotopy equivalent,'' this asks whether there are continuous maps \(u:\mathcal S\to\overline{\mathcal M}\) and \(v:\overline{\mathcal M}\to\mathcal S\) satisfying
\[
v\circ u\simeq\operatorname{id}_{\mathcal S},\qquad
u\circ v\simeq\operatorname{id}_{\overline{\mathcal M}}.
\]
Here \(\mathcal S\) and \(\overline{\mathcal M}\) are placeholders: the record does not define the Sullivan space, its topology, the moduli problem, or the allowed compactification. The interpretation as homotopy equivalence also requires confirmation.
\item Is \(T^\varepsilon(S_g)\) a Stein manifold admitting a plurisubharmonic exhaustion? For a complex manifold \(U\), a continuous function \(\rho:U\to\R\) is plurisubharmonic if its pullback along every holomorphic disk is subharmonic; it is an exhaustion if \(\{x\in U:\rho(x)\leq c\}\) is compact for every \(c\in\R\). The source does not specify additional smoothness or strict plurisubharmonicity. The set \(T^{>\varepsilon}(S_g)\) is an open subset of \(T(S_g)\) and inherits a complex structure. The closed inequality does not by itself provide a complex manifold without boundary, so its interpretation in the Stein question must be specified. The record also does not state whether the question concerns every admissible \(\varepsilon\), sufficiently small \(\varepsilon\), or some \(\varepsilon\).
\item What is the Dolbeault cohomology of the intended thick part? For a specified complex manifold \(U\), write \(\mathcal A^{p,q}(U)\) for smooth complex differential forms of type \((p,q)\), and define
\[
H_{\bar\partial}^{p,q}(U)=
\frac{\ker\bigl(\bar\partial:\mathcal A^{p,q}(U)\to\mathcal A^{p,q+1}(U)\bigr)}
{\operatorname{im}\bigl(\bar\partial:\mathcal A^{p,q-1}(U)\to\mathcal A^{p,q}(U)\bigr)},
\qquad 0\leq p,q\leq\dim_{\C}U,
\]
with zero denominator when \(q=0\). This gives the usual smooth-form interpretation if the intended thick part is an open complex manifold. For a thick part including its boundary, the function spaces, boundary regularity, and any boundary conditions must instead be specified. The record supplies none of these choices.
\end{enumerate}
\subsection{Short English statement}
The source asks which Deligne--Mumford classes arise from a Penner compactification, whether a Sullivan space has the homotopy type of a suitably compactified moduli space, and whether a thick Teichmüller part is Stein with a plurisubharmonic exhaustion, together with a computation of its Dolbeault cohomology. The compactifications, comparison maps, parameter ranges, and treatment of the thick-part boundary require clarification. The supplied open-status label and literature assessment are dated curation, not independently verified conclusions.
\subsection{Sources}
\begin{itemize}
\item Supplied local UnsolvedMath registry record, permanent ID \(30001639\), source problem code \texttt{catalogue entry}, title \emph{Compactifications and Stein Structures of Thick Teichmüller Parts}, in the Oberwolfach Reports Open Problems collection. No individual proposer is identified.
\item Supplied report citation: \emph{Teichmüller Theory (2011)}, Oberwolfach Reports, DOI \texttt{10.4171/owr/2010/53}, pp.\ 3085--3157. Original Oberwolfach problem report: \url{https://doi.org/10.4171/owr/2010/53}. The record gives proposed year 2010, while its report title gives 2011; both are retained without resolving their relationship.
\item Supplied literature source labeled ``deformation retract to the thick part'': \url{https://arxiv.org/abs/1302.0877}. The record supplies no authors, precise theorem, or parameter hypotheses for this reference.
\item Literature checked at \texttt{2026-08-21}: the supplied assessment labels the problem open and states that the thick part is contractible and has deformation-retraction results, while the Stein and Dolbeault-cohomology questions remain open in the literature checked. This is supplied dated curation, not independently verified; the background explicitly requests independent verification of time-sensitive status and cited claims.
\item The supplied registry labels the statement \texttt{corrected\_verified} and reports that its mathematical content was checked against the official Oberwolfach report, with notation, context, and question form normalized. This is a supplied verification claim; no independent source consultation is asserted here.
\end{itemize}
\end{document}
